% Generated draft: requires mathematical and source review.
\documentclass[11pt]{article}
\usepackage[margin=25mm]{geometry}
\usepackage{fontspec}
\setmainfont{Latin Modern Roman}
\usepackage{amsmath,amssymb,mathtools,unicode-math}
\setmathfont{Latin Modern Math}
\usepackage{enumitem,needspace,xurl,hyperref}
\setlength{\parindent}{0pt}
\setlength{\parskip}{5pt}
\setlength{\emergencystretch}{3em}
\newcommand{\N}{\mathbb{N}}
\newcommand{\Z}{\mathbb{Z}}
\newcommand{\Q}{\mathbb{Q}}
\newcommand{\R}{\mathbb{R}}
\newcommand{\C}{\mathbb{C}}
\newcommand{\F}{\mathbb{F}}
\newcommand{\A}{\mathbb{A}}
\newcommand{\PP}{\mathbb{P}}
\newcommand{\E}{\mathbb{E}}
\newcommand{\eps}{\varepsilon}
\newcommand{\dd}{\,\mathrm{d}}
\begin{document}
\section*{UnsolvedMath ID: 10900040}
\textbf{Source problem code:} catalogue entry
\subsection{Definitions and mathematical statement}
The source fixes a hyperbolic three-manifold \(M\) and asks, for every positive integer \(n\), for an ideal triangulation with \(m\geq n\) positively oriented tetrahedra. For an oriented hyperbolic realization, an ideal tetrahedron is positively oriented when its signed hyperbolic volume is positive; in an orientation-compatible shape convention this is \(\operatorname{Im}z>0\). This concerns the tetrahedron's geometric realization, not merely a combinatorial vertex ordering.

Let \(N(\mathcal T)\) be the total number of tetrahedra in a finite ideal triangulation, and \(N_+(\mathcal T)\) the number positive in its realization compatible with the given hyperbolic structure. The wording admits two distinct precise readings:
\[\forall n\geq1\ \exists\mathcal T\text{ of }M:\quad N_+(\mathcal T)\geq n,\]
or the stronger all-positive reading
\[\forall n\geq1\ \exists\mathcal T\text{ of }M:\quad
N(\mathcal T)=N_+(\mathcal T)\geq n.\]
In the first, other tetrahedra may be negative or degenerate if the admissible realization allows them; in the second every tetrahedron is positive. The supplied excerpt does not explicitly state which reading its \(m\) denotes, and the surrounding research summary discusses both positively oriented triangulations and numbers of positive tetrahedra. No silent choice between these targets is made.

The scope of \(M\) is also essential. In the usual topological definition, an ideal triangulation removes the vertices from a finite tetrahedral face-pairing space and identifies the resulting interior with \(M\); the familiar complete finite-volume setting is cusped. The question says only hyperbolic, without orientability, completeness, finite volume or cusp assumptions, and does not define ideal triangulations for closed cases or specify the required geometric gluing/straightening data. These conventions must be fixed before the displayed assertions can be decided. All triangulations sought are of the given \(M\), not merely of its finite covers.
\subsection{Short English statement}
For every threshold, seek an ideal triangulation of the same hyperbolic manifold with at least that many positive tetrahedra. It is unclear whether every tetrahedron must be positive, and the manifold and realization conventions are unstated.
\subsection{Sources}
\begin{itemize}
\item Local UnsolvedMath registry, permanent ID \(10900040\), exact source code \texttt{catalogue entry}.
\item Gabai and Trnkova, attribution in Delp, Hoffoss and Manning, \emph{Problems In Groups, Geometry, and Three-Manifolds} (2015), subsection 5.15, PDF page 5: \url{https://arxiv.org/pdf/1512.04620#page=5}; abstract \url{https://arxiv.org/abs/1512.04620}. Registry access 2026-07-29; source-tex extraction recorded as checked against rendered PDF text.
\item The summary mentions Gu\'eritaud, Choi and others and a work titled \emph{Positively oriented ideal triangulations on hyperbolic three-manifolds}, without full bibliography.
\item \emph{Infinitely many virtual geometric triangulations}, \url{https://arxiv.org/abs/2102.12524}, as supplied related literature; authors are not given in the record.
\end{itemize}
\end{document}
